\documentclass[12pt]{article}
\usepackage{xeCJK}  

\usepackage{charter}
\usepackage{fullpage}
\usepackage[colorlinks=false]{hyperref}
\usepackage{ifthen}
\usepackage{comment}
\usepackage[title,titletoc]{appendix}
\usepackage{pagecolor}
\usepackage{amsmath}
\usepackage{amsfonts}
%\usepackage[normalem]{ulem}
\usepackage{siunitx}
\usepackage{amsthm}
\sisetup{per=slash, load=abbr}

\usepackage{pgfplots}
\usetikzlibrary{positioning}
\usetikzlibrary{fit}
\usetikzlibrary{snakes}
\usetikzlibrary{shapes.geometric}
\usetikzlibrary{patterns}
\usetikzlibrary{shapes,arrows,chains}
\usepgfplotslibrary{patchplots,colormaps}
\usetikzlibrary{calc}
\usetikzlibrary{positioning, fit}
\usetikzlibrary{backgrounds}
\usetikzlibrary{intersections}

\newcommand{\whitepaper}[1]{\begin{center}\fbox{\parbox{0.75\textwidth}{{\small
#1}}}\end{center}}

\newcommand{\pcolor}{orange!25}

\usepackage{setspace}
\usepackage{algorithm2e}
\bibliographystyle{ieeetr}

\usepackage{geometry}
\geometry{left=3cm,right=3cm,top=1.6cm,bottom=3cm,headheight=0pt,headsep=1.5em}
\usepackage{fancyhdr}
\pagestyle{fancy}
\chead{\includegraphics[scale=0.2]{../common/Nebulas.png}}  %在此处插入logo.pdf图片 图片靠左
\lhead{} % 页眉中间位置内容
\rhead{}
%\setlength{\topskip}{1em}

\usepackage{indentfirst}


\setCJKmainfont[BoldFont = STSongti-SC-Bold]{STSongti-SC-Regular}
\setCJKfamilyfont{hei}{SIL-Hei-Med-Jian}		%宋体
\setCJKfamilyfont{song}{SimSun}		%宋体
\setCJKfamilyfont{kai}{Kaiti}		%楷体
\setCJKfamilyfont{fang}{song}	%仿宋
\setCJKfamilyfont{li}{song}			%隶书
\setCJKfamilyfont{you}{Yuanti}		%幼圆

\newcommand{\song}{\CJKfamily{song}}	%宋体
\newcommand{\hei}{\CJKfamily{hei}}	%黑体
\newcommand{\kai}{\CJKfamily{kai}}	%楷体
\newcommand{\fs}{\CJKfamily{fang}}	%仿宋
\newcommand{\li}{\CJKfamily{li}}		%隶书
\newcommand{\you}{\CJKfamily{you}}	%幼圆
\newcommand{\reffig}[1]{图\ref{#1}}
\newcommand{\refsec}[1]{\S \ref{#1}}

\onehalfspacing   % ----------设置1.5倍行距（可能有意义，待调整）

%\parindent=20pt  % -------------------首行缩进大小，英文分段就直接0pt了吧。
\setlength{\parindent}{2.1em}
\setlength{\parskip}{0.3\baselineskip}
\newcommand{\nrcore}{Core Nebulas Rank}
\newcommand{\nrext}{Extended Nebulas Ranks}
\newcommand{\dom}{{\; \texttt{dom}\;}}
\newcommand{\xpcomment}[1]{{\color{blue} #1}}

\setCJKsansfont[BoldFont = STHeitiSC-Medium]{STHeitiSC-Light}


\newtheorem{property}{特征}
%\addbibresource{reference.bib}

\begin{document}
\pagestyle{empty}
\renewcommand{\contentsname}{目录}
\renewcommand{\abstractname}{摘要}
\renewcommand{\refname}{参考文献}
%\renewcommand{\nomname}{术语表（按首字母排序）}
\renewcommand{\figurename}{图}
\renewcommand{\tablename}{表}
\renewcommand{\baselinestretch}{1.5}
\renewcommand{\appendixname}{附录}
\renewcommand{\proofname}{证明}

\pagecolor{\pcolor}

\begin{titlepage}
  \begin{center}
    \vspace*{5.5cm}
    \includegraphics[scale=0.5]{../common/Nebulas.png}
    \vspace{0.5cm}


    \textbf{\huge{星云治理橙皮书}}

    \vspace{0.5cm}
    吴涤、李晨、范学鹏、王冠、曾驭龙、刘歆子建
    \vfill
    2019年4月 \\
    版本号:0.0.9
    \textbf{}
  \end{center}

\end{titlepage}
\setcounter{page}{0}
%\thispagestyle{empty}
\tableofcontents
\newpage
\setcounter{page}{1}
\pagestyle{fancy}	
\vspace*{0.01cm}
% !TEX root = main.tex
\section{概述}

\textbf{关键字： 治理\ 去中心化\ 协作\ 激励\ 自进化\ 自治 }

\vspace{2em}

星云链（Nebulas）是开源公链，星云是自治元网络（Autonomous Metanet）~\cite{AutonomousMetanet}，星云专注于处理复杂数据和交互、复杂的协作关系，致力于通过区块链等技术手段，实现让每个人从去中心化协作中公平获益的愿景~\cite{vision}。

本橙皮书着重介绍为了实现愿景，星云如何借助于独有的创新技术\textbf{治理链上公共资产，探索全新的去中心化协作模式，实现提供正向激励、可自进化的去中心化自治组织}（Decentralized Autonomous Organization，DAO）~\cite{DAO}。星云治理在如下三方面进行了独创设计，本书将一一介绍：

\begin{enumerate}
	\item \textbf{组织结构和监督机制：}
	独立运行且互相制约的星云三会：星云理事会、星云基金会、星云技术委员会，阐述其基本构成、权力和义务；
	\item \textbf{链上协作方式：}
	星云社区项目制介绍，星云通过“NAT链上投票”（参见\ref{nat}）实现社区协作和系统升级的流程；
	\item \textbf{经济体和激励：}
	在星云治理过程中，星云经济体如何对每一位社区成员提供正向激励。
\end{enumerate}
\section{Background}

The goal of Nebulas governance is to achieve the Nebulas vision and it focuses on decentralized collaboraation. Before introduce the details of Nebulas governance, you need to know the difficulty of collaboration in the new situation now and which problems that Nebulas wants to solve.

\label{background}

Human beings are social aand we are not strangers ot Collaboration. Even Robinson on the island had a set of models to collaborate with Friday~\cite{robinson}. The collaboration itself does not have absolute advantages and disadvantages. In different situations, one or more than one collaborations methods suitable for the situations should be selectd. With the developmeent of science and technology, the situation has been developed from face-to-face cooperation to a global collaboration accross the locations and organizations. The goal of collaboration is also more variable, the results raanging from physical to virtual, and the period becomes longer and more flexible.

Nebulas doesn't want to subvert other forms of collaboration, nor does it exclude multiple ways of collaborating. Nebulas tries to find a more appropriate way to collaborate to complement others in new situations. The new situations has the following features:

\begin{itemize}
	\item \textbf{Information interaction is developing from simple to complex.}

	The birth of the decentralized digital cryptocurrency, Bitcoin, allows the blockchain to records transaction information. The second-generation blockchain system, such as Ethereum, proposes a smart contract with Turing completeness, and the blockchain becomes programmable. Then there are more and more data and assets iinteraction problems in different situations such as on-chain, off-chain, and cross-chain.

	\item \textbf{The user roles are also increasing.}

	In the early Bitcoin community, there were only miners and token-holders. After Ethereum, there were developers, users. And more and more people know blockchain, then it becomes a challenge to distribute of rights and responsibilities of different user roles.

\end{itemize}

The problems such as:

\begin{enumerate}
	\item 

	\textbf{Centralized governance can't face to the complex new situation.}

	Blockchain is essentially a decentralized, untrusted, game-based autonomous system. Its true charm is the open collaboration model based on the consensus mechanism and the decentralization idea~\cite{whitepaper}. However, some blockchain projects simply do the opposite and use centralization way to govern, such as directly arbitraating hackers through the core arbitration court. The legitimacy and fairness of this approach are difficult to guarantee. Faced with complex data interaction patterns and rich user roles, the centralization of single evaluation criteria is difficult to achieve a wide and comprehensive, causing community members to resist. On January 11, 2019, the EOS Authority initiated a vote on whether to delete the ECAF (The EOSIO Core Arbitration Forum), and the percentage of support was as high as 98\%~\cite{DeleteECAF}。

	\item 

	\textbf{The existing decentralized governance rules are not uniform.}

	For example, in the Bitcoin community, for users of different roles such as miners and holders, the corresponding rules are not the same and are not clear. Such decentralized governance methods are likely to cause unclear community development goals, and it is difficult to effectively organize and iteratively upgrade.

	\item 

	\textbf{Traditional decentralized collaboration is often in the tragedy of the commons~\cite{TragedyOfTheCommons}.}

	Traditional decentralized collaborative projects, such as a large number of open source communities, have unclear interest models. Funding sources are often based on donations. Upgrades and evolutions rely too much on developers' interests. Frequently, there are problems such as the tragedy of the commons and the slow evolution of ecology. There are more people using public resources (such as open source code) and fewer people contribute. Open source projects that rely on large companies and enterprises to donate donations are often blocked by the development direction of large companies and become affiliates of enterprises.

	Blockchain technology, because of the existence of tokens, gives us the opportunity to solve the basic dilemma of decentralized collaboration by providing sustainable incentives and building a benign economy.

	\item

	\textbf{The early blockchain project consensus mechanism incentives are not comprehensive, and community participation is low. }

	For example, the Proof of Work (PoW~\cite{pow}) used by Bitcoin only focuses on mining incentives. This single incentive cannot cope with the gradual enrichment of user roles. Decentralized Ethereum has been criticized for its slow upgrades. Ethereum's upgrade proposal needs to be widely recognized by the community and then executed by the miners. However, the opinions of the parties in the entire ecology are difficult to unify. The incentives of Ethereum do not cover the different user roles of the entire ecosystem. The phenomenon of “doing nothing” has led to low participation in the Ethereum upgrade proposal, delay in implementation, and impediment to ecological development.

\end{enumerate}

There is currently no perfect solution to solve the above problems. We realize that in the new world with unprecedented complexity, the birth of new technologies is expected.


\section{星云治理概述}
星云致力于通过技术手段，减小人为治理通常存在的摩擦和治理成本，用科技改变协作关系，促进社区良性发展，实现星云愿景，让每个人从去中心化协作中公平获益。

\subsection{基本主张}
一切复杂系统的演绎从制定符合“常识”的基本规律开始。区块链资产的基本构成单位是“地址”，也是星云社区治理的基本单位。在此，我们正式提出关于星云链“地址”三点基本权利主张：
\begin{enumerate}
	\item 拥有和操作星云链上资产的权利；
	\item 发起提案的权利；
	\item 参与提案投票的权利。
\end{enumerate}
星云坚信每个地址在系统中拥有的以上基本权利，在任何条件下不受侵犯。而每个地址所对应的私钥是该地址控制权的唯一凭证，通过它，每一位星云社区成员都享有使用星云主网、参与社区决策、参与生态建设的权利。没有绝对中心化组织或个人，任何社区成员都可以提出提案，通过社区链上投票进行治理。星云未来发展方向的决定权在每一位星云社区成员手中。

\subsection{治理范围}
区块链是承载协作关系的网络，亦是承载激励协作关系“资产”的网络，在没有绝对中心化权力的网络里，相应的公共资产理应属于全体社区成员共同管理。

星云治理所针对的公共资产主要包括：
\begin{enumerate}
	\item 公共开源代码知识产权（如主网升级等影响到星云链上公共利益的相关代码）；
	\item 白皮书公示的预留公共资金池。
\end{enumerate}


\section{治理机构}
现阶段星云治理机构包括三部分：
\begin{enumerate}
	\item 星云理事会
	\item 星云基金会
	\item 星云技术委员会
\end{enumerate}
三个机构均允许社区所有人自愿报名加入。为了保障公正透明，有2个准则：
\begin{enumerate}
	\item 每个人最多同时加入两个机构。不可在任意时间段内同时在三个机构中担任职务。
	\item 三个机构互相制约，任何机构不具有独立决策和使用公共资金的权利。
\end{enumerate}
\subsection{星云理事会}
\subsubsection{权利和义务}
\textbf{星云理事会职能}

星云理事会职能仅限于保障和监督各项提案正常推进，推动星云愿景实现。不具备公共财产的处置权。

\textbf{理事权利}
\begin{enumerate}
	\item 提请提案”二次投票“的权利，7人共同商议决策，各自表决。
	\item 以星云理事会理事的身份进行推广。
	\item 委任星云技术委员会等组织机构或个人处理星云社区公共事务。
\end{enumerate}
\textbf{理事义务}
\begin{enumerate}
	\item 参与治理流程监督。
	\item 监督社区预留资金等公共财产安全。
	\item 宣传·推广星云理念。
\end{enumerate}
\subsubsection{组织结构}
\textbf{席位}

一共7名理事。

2019年第一届星云理事会理事，由星云基金会提名3席，社区公开选举产生4席。

星云基金会提名席位每两年至少减少一席。即最晚6年后星云基金会不再提名。

\textbf{任期}

2年/期，可连任1次

\subsubsection{选举}
\textbf{选举原则}

星云理事会理事通过链上治理流程投票产生。
\begin{enumerate}
	\item 只要拥有至少一个星云主网NAS账户地址的社区成员，都有选举与被选举权。
	\item 首期星云理事会选举方案由星云技术委员会代为提出。未来方案升级需经过链上公开投票通过方可执行。
	\item 社区成员对星云理事会有监督。
\end{enumerate}
\textbf{选举流程}

提名：公开阐述自己的贡献、治理理念和方案等。时长为一个自然月。

动员：候选人可在星云社区自行组织动员\&拜票。时长为一个自然月。

选举：公开链上投票。每人可选多人。时长不超过一周。

\textbf{获选条件}
\begin{enumerate}
	\item 需要获得所有投票参与者中，星云指数（Nebulas Rank，NR，核心星云指数考虑资产中值和出入度衡量，用于衡量一段时间内用户对整个经济体的贡献）不小于51\%的支持率，详见5.2投票机制。
	\item 取排名靠前者
\end{enumerate}
\textbf{中期投票和罢免}

任期满一年的理事需要做述职。社区根据述职情况进行中期投票，如果获得NR占比不小于51\%的支持率则继续任职，整个星云理事会需要重新选举。

如果遇到此种星云理事会集体罢免的情况，由星云技术委员会组织并监督新一轮星云理事会选举正常进行。在此期间，由述职通过的理事暂时代理星云理事会日常事务，但授权范围不扩大。

\subsubsection{资产监督}
\textbf{资产来源}

星云理事会不具有社区公共资产。仅负责监督社区公共资产的使用。社区公共资产主要包括2部分来源：
\begin{itemize}
	\item 35,000,000 NAS（35\%）：社区预留

	（自星云创始起锁仓，解锁日期不早于2019年6月23日）
    \item 8,219.1744 NAS/日：共识记账系统增发
    
    （目前增发量为3\%，每日自动产生）
\end{itemize}
依照《星云非技术白皮书》（2017）所示，在2019年9月节点实现去中心化之后，共识记账系统增发将变为4\%，包括三部分：
\begin{itemize}
	\item 1\%：开发者激励协议原生激励（Developer Incentive Protocol，DIP），当前约为每周300 NAS
	\item 2\%：共识记账收入
	\item 1\%：理事会项目发展资金储备
\end{itemize}
在2019年9月节点实现去中心化之前，已经增发的部分由星云理事会负责管理，DIP部分暂由星云理事会管理。未来将要从中扣除已产生的服务器费用等相应记账开支（暂由星云基金会垫付，资产移交时进行扣除）。

\textbf{监督方案}

不同的资产用途有不同的监督方案。
\begin{enumerate}
	\item \textbf{项目协作平台}
	
	\textbf{方式:} 由星云社区全体参与成员公开审核，自动分发。
	
	社区自由提交资金申请，通过公开投票获得使用批准并且没有其他违规现象（参见2.2治理原则），即可申请到社区预留资金，在项目完成测试并上线后自动获得资金。
	
	\textbf{星云理事会职能：}如果某个项目存在违背治理基本准则的情况，星云理事会具有一次驳回重新投票再审的权利。
	
	\item \textbf{开发者激励协议原生激励}
	
	\textbf{方式：}自动，每周一次。
	
	星云是首个具有原生激励的公链。开发者激励协议（DIP）直接提供原生激励，每周一次自动向活跃DApp开发者的地址分发奖励NAS。
	
	\textbf{星云理事会职能：}链上自动化治理机制，星云理事会不具处置权，仅负责监督其安全性和稳定性。
	
	\item \textbf{星云基金会日常运作}
	
	包括星云基金会理事收入、紧急预案处理等无法计入项目制平台的事宜。此部分应有公示。
\end{enumerate}

\subsubsection{收益}
\textbf{总收入}

10,000 NAS/2年

\textbf{收益发放}

每满半年为一期，2年内分4次分发。每次分发额度为：3,000 NAS、4,000 NAS、6,000 NAS、7,000 NAS。如中期投票没有通过，后两期收入不予发放。

\textbf{财务要求}

为保证经济体利益的一致性，和星云理事会政策的延续性，理事正式入职时需要抵押100,000 NAS。卸任半年后解锁返还。

\subsubsection{公开透明}
星云理事会应确保治理流程和资产使用公开透明。
\begin{enumerate}
	\item 通过组织撰写季报等手段，定期公示资产和社区发展等情况。
	\item 如产生技术升级、项目申请驳回重新投票等处理情况，应及时公示。
	\item 所有人员选举情况、委任情况（如委任技术委员会管理项目协作平台）等，应及时公示。
\end{enumerate}

\subsection{星云基金会}
星云团队于2017年6月组建，星云基金会随之成立。负责星云团队的运维，保障项目正常进行。在实现《星云非技术白皮书》（2017）中承诺的技术点之前，星云基金会将继续肩负其历史使命。
\subsubsection{权利和义务}
\textbf{理事权利}
\begin{enumerate}
	\item 基金会内部选举和被选举权。
	\item 享受星云基金会运营带来的收益。
	\item 参与基金会发展和投资等的决策。
\end{enumerate}
\textbf{理事义务}
\begin{enumerate}
	\item 管理星云非公共预留部分的资产，包括星云团队成员奖励发放等。
	\item 按照星云发展需求组织生产，保障星云项目的研发可以正常进行，按时完成《星云非技术白皮书》（2017）承诺的开发内容和时间点。
	\item 每年一次对星云理事会述职，持续为星云生态服务。
\end{enumerate}
\subsubsection{组织结构}
\textbf{席位}

常务理事不小于5席，其中，主席1席，秘书长1席。

理事若干。

\textbf{任期}

1年/期，可连任

\subsubsection{选举}
\textbf{入选资格}

在星云基金会管理的团队资金池中，享有奖励额度达到80,000 NAS后即自动拥有进入星云基金会的资格。

\textbf{选举}

主席由现任星云基金会理事内部推选，支持率需要过半数。每位理事具有被选举权和选举权，一人一票。
由新上任的主席委任秘书长和其余常务理事，不少于3名。

\textbf{弃权}

每位具有加入星云基金会资格的人员，均有放弃理事资格的权利。
当常务理事不足5名时，则依照奖励额度排位候补。

\textbf{罢免}

星云基金会通过内部决议具有罢免星云基金会理事的权利。但结果必须向社区公示。
被罢免的星云基金会理事有权利向社区公开述职和要求通过全员投票重新裁决。

\subsubsection{资产管理}

\textbf{资产来源}

星云基金会管理的资产包括：
\begin{itemize}
	\item 20,000,000 NAS（20\%）：星云团队预留
\item  5,000,000 NAS（5\%）：星云社区发展基金剩余部分
\item 早期私募所得项目发展资金
\item 早期生态投资所得
\end{itemize}

\textbf{资产管理方案}

星云基金会资产发放采用多重签名的方案：5签3～5签5，则资金可用。（参考NAS当前流动性和总额进行加权，涉及使用的金额越大，需要的签名就越多。）

\subsubsection{收益}
\textbf{总收入}

\begin{enumerate}
	\item 星云基金会薪资
\item 工作每满1年，如述职通过，即可享受NAS奖励（由星云基金会常务理事负责分配）
\item 享受星云基金会生态投资等带来的收益
\end{enumerate}

\textbf{财务要求}

为保证经济体利益的一致性，和星云基金会政策的延续性，理事正式入职时需要抵押50,000 NAS。卸任半年后解锁返还。

\subsubsection{公开透明}
星云基金会应确保资产整体使用情况公开透明。
\begin{enumerate}
	\item 通过半年报等形式，披露整体经济情况和支出情况。
	\item 公示星云团队预留部分（20\%）的整体分发情况。
	\item 大于500,000 NAS的大额支出，需要向社区提前披露。 
\end{enumerate}

\subsection{星云技术委员会}
星云技术委员会成立于2018年9月。成立至今，星云技术委员会秉持着开放、共享、透明的精神，致力于推动星云链技术研发逐步去中心化、社区化。区块链技术为建设全新的、自激励的开源社区带来可能性。星云链的价值发现体系、自进化和自激励的技术理念为构建去中心化协作的世界提供了保障。星云技术委员会将全力推动星云愿景的实现。

自星云理事会成立后，原先由星云团队核心成员构成的星云技术委员会将完成历史使命，转型为社区化团队。星云技术委员会将受星云理事会委托，负责一系列社区化相关平台的日常运作，以及流程正当性监督。

\subsubsection{过渡期}
在2019年星云理事会第一次选举完成之前，星云技术委员会暂由星云团队成员代理。在原先星云技术委员会成员基础上进行拓展。星云理事会第一届理事人员确定后，星云技术委员会将启动改组。

\textbf{过渡期委员名单}

徐义吉、王冠、范学鹏、陈聪明、王卓尔、李晨、吴涤、温晨阳、鲁妍芬

\textbf{过渡期职责}
\begin{enumerate}
	\item 监督星云项目发展计划顺利进行
	\item 项目制平台和投票机制等社区化相关工具上线
	\item 第一次全网链上投票升级IR
	\item 监督星云理事会选举流程
\end{enumerate}

\subsubsection{权利和义务}
\textbf{委员权利}
\begin{enumerate}
	\item 参与项目审核和测试可以按照所审核项目价值提取项目申请额度的5\%作为奖励回馈，多劳多得。
	\item 享有星云技术专家团队的荣誉。
\end{enumerate}

\textbf{委员义务}

\begin{enumerate}
	\item 社区化提案质量监督。如果提案遭到举报，或者发现异常提案，三人核实后向理事会发起重审。
	\item 技术类提案的测试。技术把关。
\end{enumerate}

\subsubsection{组织结构}
\textbf{席位}

不限
\vskip 10pt

\textbf{任期}
1年/期，可连任

\subsubsection{选举}
自荐和社区推荐相结合，对社区公开述职。

候选人需通过技术委员会技术专家的技术能力考核，并通过全体成员投票（一人一票，支持率过半），方可决定是否加入。

\subsubsection{资产管理}
星云理事会委任技术委员会维护流程正当性，相关平台和工具正常运作。星云技术委员会本身不持有资产。

\subsubsection{收益}
\textbf{总收入}
\begin{itemize}
	\item 委托佣金（按月发放）
	\item 项目制相关审核和监督工作的5\%提成
\end{itemize}

\textbf{财务要求}

为保证经济体利益的一致性，和星云技术委员会政策的延续性，理事正式入职时需要抵押25,000 NAS。卸任3个月后解锁返还。


\section{治理流程}
在NBRE中，升级通过提交IR（LLVM IR）来进行。不同IR的提交会对星云链整体系统环境的性能产生实质性的影响，IR的提交实际程度上决定了星云链主网的升级和性能推进方向。

现阶段星云功能特点未全部实现，星云团队会慎重开放提交IR的权限。当前治理流程如下：

\textbf{治理流程（现阶段）}	
\begin{enumerate}
\item \textbf{提案阶段}（Proposal Period）：在社区公开讨论提案，发起项目，并可申请使用社区预留资金；
\item \textbf{投票阶段}（Proposal Voting Period）：投票区块期间，赞同票NR占比超过全网NR值一定比例x\%，视作通过；
\item \textbf{开发阶段}（Develop Period）：提案人进行开发，或在社区寻找开发者共同开发，开发过程在协作平台（go.nebulas.io）上公开；
\item \textbf{测试阶段}（Testing Period）：经主网团队操作，部署测试网公测，关于IR是否可以部署到主网进行链上加权投票，投票区块期间，赞同票NR占比超过全网NR值不小于x\%，视作通过；
\item \textbf{发布阶段}（NBRE Period）：经技术委员会验收核实，且测试期内没有争议，IR被NBRE执行，并入主网。
\end{enumerate}

\textbf{治理流程（NF实现后）}

NF实现后，任何人可以提交IR，以主网NF侧链形式存在，不需要经过主网团队审批即可自行测试。

\subsection{提案}
星云链技术开发和运营为项目制，项目的发起、立项、预决算、执行、审核由星云技术委员会负责组织。星云技术委员会只负责监督，保证整个流程公开透明。如涉及线下活动，应符合当地相关法律法规。原则上星云技术委员会不介入其中，并且不承担任何责任。

\textbf{提案人}

所有社区成员。

\textbf{提案方式}

在提案专区公开发出提案。

\textbf{提案内容}

一份完整的提案需要包括以下内容：
\begin{enumerate}
\item 提案人姓名（需要登录）
\item 提案描述（当前问题、解决方案、理由）
\item 提案成果形式（如github地址、DApp、PDF、文稿链接等，也可以是综合性结果或报告）
\item 指定执行人，无指定执行人则可以在社区公开招募执行人
\item 各时间点设定：投票开始和结束时间、招募执行人时间、执行阶段时间点和截止日期、执行耗时预估、测试时长
\item 如果申请使用社区预留资金，需要填写明确金额（NAS作为计价单位）
\end{enumerate}

\textbf{提案时长}

由提案人自由定义。

\textbf{提案修改}

在投票正式开始之前，提案人可以修改提案所有内容。其他人可以评论、关注等。每次修改需要支付$n^2$ NAS。这部分奖励会均分给当次提案前参与过讨论的用户。

\textbf{提案取消}

在正式投票开始之前，提案人可以取消提案。正式投票开始后，不可取消提案。

\subsection{投票}
\subsubsection{星元币}
星云生态中的治理将通过社区投票表决实现，星云生态中的投票将使用星元币（Nebulas Autonomous Token，NAT）作为唯一的投票介质。

\textbf{名称}
 
星元币 Nebulas Autonomous Token

\textbf{代号}

NAT

\textbf{形式}

NRC20 token

\textbf{用途}

NAT将作为星云生态中投票场景的唯一介质，支持星云理事会、星云基金会和星云技术委员会推进星云社区的治理。社区成员可以通过使用NAT投票表达自己自己对星云生态的意见，包括但不限于星云理事会的选举、星云主网NBRE中执行的IR及星云项目提案的立项表决等。

\textbf{发行方式}

NAT将按照一定周期通过空投的方式投放给每个星云主网地址的持有者，投放的比例与该地址的在该周期之内的星云指数（Nebulas Rank，NR）及该地址的投票行为相关。

关于NAT的具体的发行方式和参数请参见【5.3】NAT算法

\textbf{发行量}
	
NAT的发行量与星云链原生代币NAS的供应量以及星云链全网的星云指数（Nebulas Rank，NR）相关。在星云链全网星云指数恒定的情况下，NAT发行量恒定且释放的数量逐年递减。按照当前星云链全网的星云指数的情况，NAT的总发行量约为【】枚，且每两年增发的数量减少【50】%。（待确认）

【插入发行量递减公式】：【】

\textbf{管理}

用户可以在NAS Nano Pro及其他支持NRC20的钱包中管理自己的NAT。同时，用户可以在支持星云链的区块链浏览器上（比如explorer.nebulas.io）像查看其他NRC20 token一样查看NAT的交易、流通情况等。

\textbf{投票交易}
用户的投票将在星云链主网上实现，用户可以将自己手里的NAT以星云链主网上NRC20代币交易的形式转账给星云链投票的DApp【插入DApp名称】，转入星云链投票DApp的NAT将被烧毁。

\textbf{使用NAT作为投票介质的理由}

星云链聚焦链上治理，致力于运用区块链技术来提供更公平的协作环境。根据星云链的治理原则，链上治理的最基本单位是一个星云主网地址，星云生态的治理投票将通过链上投票的方式进行。链上投票所有的投票行为和交易信息均将在区块链上记录，公开、透明、不可篡改。

同时，星云链期望运用区块链技术构建一个面向复杂数据和交互的，可以促进社区良性发展的协作体系。在这个体系中，对于生态具有更多贡献的人应该在生态治理中享有更多的话语权，应该被激励更多的投票权。

基于此，一段时间内某个地址的NAT的持有量应该要反映一段时间内该地址持有者在生态中的贡献度。同时，NAT需要具备良好的抗作弊性，需要能够防止恶意的刷票行为。


根据NAT的发行方式，某个地址在一段时间内获得的NAT的数量将主要受两个因素影响：1）该地址的在一段期间内的核心星云指数，2）该地址的投票行为。


其中星云指数是星云链的独创技术，核心星云指数也是将在Nebulas NOVA中重要功能之一。
根据2018年6月30日星云研究院发布的《星云指数黄皮书》，星云指数用以衡量区块链世界链上数据价值，以流动性、传播性、互操作性等体现数据互动关系的因素为基础，可以用来衡量地址、智能合约、去中心化应用等对象的影响力。其中，核心星云指数
(Core Nebulas Rank)考虑资产中值和出入度衡量，在微观层面，核心星云指数衡量了每个账户对经济系统的贡献；在宏观层面，核心星云指数则用经济学中经典的货币数量方程描述了区块链中货币数量、货币价值和流通速率以及生产力几者之间的关系。
因此，我们将每个账户可以获得的NAT的数量和一段时间内每个账户的核心星云指数相关联，并将NAT在一段时间的总供应量和整个星云生态的核心星云指数相关联。一方面，从微观层面反映每个账户在一段时间内对经济体的贡献程度，另一方面使NAT的总供应量和整个经济体宏观层面的活跃程度相关。

除此之外，公平性是星云指数的重要特性之一。在《星云指数黄皮书》我们已经指出PageRank类排名策略在应对“女巫攻击”时存在缺陷。在区块链系统上，女巫攻击表现为攻击者创建大量账号并且拆分资产分开进行交易获取更高收益。星云指数将主要通过以下两点来应对操作和作弊带来的挑战
1)	旨在通过防止用户借由控制多个账户实现某交易行为的收益大于通过单一账户实现该交易的收益
2)	星云指数中的二维出入度计算函数也鼓励账户之间的正常交易流通
因此，在星云链主网上的”女巫攻击“行为无法为攻击者带来更高的星云指数的收益，进而无法为攻击者带来超出其对经济体实际贡献的NAT，可以有效的抵抗恶意刷票。

同时，我们认为，在生态中积极参与投票的用户应该被激励。因此，系统将会按照一段投票期间内地址的投票行为给予额外比例的NAT。没有参与的用户在一个空投周期中仍可按照自己地址的NR比例获得一定数量的NAT，但是无法获得投票行为激励的部分。

基于此，NAT不仅可以反映单位地址对于星云生态的贡献度，并且具有抗作弊性，可以应对投票过程中由于恶意刷票可能产生的舞弊行为。

我们认为在当前阶段，NAT是合适的投票介质。诚然，NAT在实际应用中仍会面临不同的挑战，我们也将根据NAT在投票场景中的实际情况及社区的意愿调整NAT的参数和算法。

\subsubsection{投票方式}
\textbf{投票种类}

每人可投：赞同、反对、弃权三种票。亦可不参与投票。

\textbf{投票时长}

提案人自定义。但必须限定投票时长，不能无限投票。提案人可以向星云技术委员会咨询指导建议。

\textbf{投票管理}

公开投票由技术委员会负责组织和管理。公开投票接受所有人的公开监督。针对违反社区治理基本准则的提案，技术委员会可以向理事会发起重审提案申请。星云理事会作为星云社区治理的监管机构，有权对任何提案的投票发起二次投票。对于同一提案，理事会可以且仅可以发起一次二次投票。

\textbf{投票流程}
\begin{enumerate}

\item 用户投票

用户在提案平台可以查看，每个地址针对一个提案可以投X NAT，投出的NAT在系统中将全额烧毁。

\end{enumerate}

\textbf{投票结果}

满足以下条件算投票通过：

\begin{enumerate}
\item 赞同票总占比大于51\%
\item 投票期间，参与投票的NAT占比超过全网NAT值一定比例x\%，x\%的决定因素为提案申请的金额占链上流通量比例。{\color{red}（建议改成固定NAT数值，可动态调）}
\item 如果金额为0，x\%有保底值。
\end{enumerate}

\subsubsection{公开透明}
投票过程时刻保持公开透明可见。投票结果不可修改。并且通过技术手段保障流程正当性。

\textbf{技术合法性}
\begin{enumerate}
\item 链上投票保证了投票过程的公正性和安全性。同时若引入匿名制还能有效保护成员隐私（隐藏其投票内容）。
\item NR的抗作弊性质保证了投票过程的公平性：成员建立大量新账户不会提高他的投票效用，同时高资产成员将其资产分散到多个子账户亦不会提升其投票总效用。
\item 投票所采用的各项技术是与Nebulas NOVA的技术一脉相承的，如DIP本质上实现的就是用户给DApp投票的过程。同时我们今后也会不断对投票技术进行补充和完善，将其作为一个不可分割的整体。
\end{enumerate}

\subsection{NAT算法}
\textbf{空投部分}

我们分若干个投票周期。我们用$n$表示当前投票周期的序号。如前所述，每个投票周期（一般而言为一周）根据每个地址的NR值获取增发的NAT。若某地址的NR值为$x$，则第$n$周期增发的NAT为$\lambda f_n(x)$。关于$f_n(x)$的选取，现给出如下方案：

%其中$\lambda \leq 1$为衰减比例参数，随当前市场现存总NAT而变化（总NAT越多$\lambda$越小）。$f(\cdot)$为NR到NAT的映射函数，递增，且满足$f(0)=0$。

\begin{itemize}
	\item $f_n(x)=g(x)\lambda^n$, $\lambda<1$。其中$g(x)$为从NR到NAT数量的一个线性映射。该方案可以理解为，增发的NAT总量每轮缩减$(1-\lambda)*100\%$
	\item 增发总量根据当前现存NAT总量动态调整。如根据上面算法设定一个等比数列预期值，然后补上预期值与当前差值。
\end{itemize}

%因为NR总量不超过NAS总量，NAS总量根据增发规律少于$10^9(1+0.002n)$，故可得$x<10^9+0.002n$。据此可证方案123 NAT总量均有上限。（具体证明通过后给出）

\textbf{投票部分}

每个用户支付的用于投票的NAT立即销毁。

\textbf{返还部分}

一周期投票结束后，为了激励正常投票的用户，会根据用户此周期内所投票数返还一定比例的NAT。返还比例为$\mu>1$乘以投票数目。最高不超过$\mu C_{in}$。其中$C_{in}$为预期用户$i$在第$n$周期投出的NAT数目。一般而言选取$C_{in}=f_n(x_i)$，其中$x_i$为此用户该周期的NR，即，我们希望用户把此周期增发的NAT用于投票。用户一次投出过多的票相当于该用户要“动老本”进行投票。

注意:$C_{in}$一定要和增发数量挂钩，以防止滚雪球效应。	

\begin{comment}
具体数额为
$$\mu C \min \{\frac{C_0}{C},1\}$$
其中$\mu>1$为激励系数，可选取$\mu=1.1$。$C$为用户此阶段（一个增发周期内）投票花费（即烧毁的）NAT总数目，$C_0$为预期希望用户投出的NAT数目，每个用户对应不同的$C_0$，一般而言选取$C_0=\lambda f(x_1)$，即我们希望用户把此阶段增发的NAT用于投票，而不推荐用户一次投出过多的票。

对于用户而言，假设他一个阶段内用$C$枚NAT进行投票且$C \leq C_0$，则其NAT的变化量为
$$\mu C \min \{\frac{C_0}{C},1\}-C = \mu C_0-C$$
上式随$C$递减。这意味着对于一个利益最大化的用户而言他的最优策略为投出$C_0$票。如若项目对其意义重大，需要投出大量NAT，则有可能$\mu C_0-C <0$，相当于该用户要“动老本”进行投票。

若$C < C_0$，则其NAT的变化量为

$$\mu C \min \{\frac{C_0}{C},1\}-C = (\mu-1)C$$
上式随$C$递增，这意味着最优策略同样是投出$C_0$票。当投票数目较少时仍能获得收益但相对投$C_0$而言较少。

\end{comment}
[通过后，此处加变化曲线图]

注意：
\begin{itemize}
\item 目前版本暂定赞同票与反对票没有区别，即返还比例不同。之后可根据票种设定并乘上不同的返回参数$\mu_1$
\item 若考虑到投票完成后系统的总NR变化，则可再乘上一个系数$\mu_2$，用于反应该周期系统的繁荣度。
\item 若上述两条采用最终返回NAT数值为
$$ \mu\mu_1\mu_2 C \min \{\frac{C_0}{C},1\}$$
\end{itemize}


\begin{property}
本算法能满足NAT总量的收敛性，即NAT总量在任何时候都不会超过一个上限。
\end{property}
\begin{proof}
	根据星云链技术白皮书的设定\footnote{https://nebulas.io/docs/NebulasTechnicalWhitepaperZh.pdf}，NAS的固定总量为$10^9$，平均每周大约增发（在固定总量的基础上）$0.2\%$，故在第$n$个投票周期市面上现存NAS总量不会超过$10^9(1+0.002n)$。
	
	接下来我们证明所有地址一个周期内的资产中值（见星云指数黄皮书中的定义\footnote{https://nebulas.io/docs/NebulasYellowpaperZh.pdf}）总和不会超过市面上现存NAS总量。这是因为，对于任意一笔数量为$y$的NAS资产，他只能最多在一个地址内存在该周期一半以上的时间（三天半），故最多给全网节点的总资产中值提供$y$的贡献。
	
	同样根据星云指数黄皮书的设定，任何一个地址的NR值不会超过该地址的资产中值（指同样一个周期内，注意NR，NAT的计算都是以周为单位，具有同步性），这是因为黄皮书NR计算公式$\Omega(\cdot)\Psi(\cdot)$中，以资产中值为输入的Wilbur函数$\Omega(\cdot)$满足$\Omega(x)\leq x$，且出入度函数$\Psi(\cdot)$值域不超过1。
	
	结合上述结论，可得在第$n$个周期，所有地址NR总和不超过$10^9(1+0.002n)$，从而NAT增发数量不超过$g(10^9(1+0.002n))\lambda^n$
	
	又因为投票返还部分的NAT不超过增发部分乘以$\mu$，故即使加上返还部分带来的增量，周期$n$内NAT的总增量超过$\mu g(10^9(1+0.002n))\lambda^n$
	
	最后，欲证NAT总量的收敛性，只需证明级数
	$$\sum_{n=1}^{\infty} \mu g(10^9(1+0.002n))\lambda^n$$
	收敛，由于$g(\cdot)$为线性函数，故
	$$ \lim_{n\rightarrow \infty} \frac{\mu g(10^9(1+0.002(n+1)))\lambda^{n+1}}{\mu g(10^9(1+0.002n))\lambda^n} = \lambda <1$$
	由比试判别法可得该级数收敛，证毕。
\end{proof}
同时，上述投票算法具有下列良好性质。
\begin{enumerate}
	\item 抗滚雪球效应：如若简单的按固定比例返还NAT，则一个用户可以每次投出所有的NAT并享受大于1比例的返还（如1.1）,则其总NAT将按$1.1^n$指数级上升，增长过于庞大。
	\item 抗收买性：若一个低NR的用户以购买的方式获取大量NAT并用于投票，由于对低NR用户我们设定的对应$C_0$较低，返还的NAT很少，大部分都被烧毁，导致该用户剩余NAT很少作为惩罚。
	\item 抗通货膨胀：由于系统增发NAT比例与当前市场NAT总量有关，可有效控制NAT的贬值。
	\item 头部效应：早期拥有高NR的用户能拥有更高NAT总量
\end{enumerate}

\subsection{执行}
\textbf{提案人义务}
\begin{enumerate}
\item 监督执行人执行，各时间点没有错失。
\item 维护社区讨论。
\end{enumerate}

\textbf{提案人权利}
\begin{enumerate}
\item 享有提案中涉及的应得回报。
\item 在申请参与执行的社区用户中挑选合适的人来执行。
\end{enumerate}

\textbf{执行人义务}
\begin{enumerate}
\item 定期公示进度，保证执行过程公开透明。
\item 保证按时保质完成，并按照提案要求提交成果。
\end{enumerate}

\textbf{执行人权利}
\begin{enumerate}
\item 享有提案中涉及的应得回报。
\item 在不影响进度并且能达成最终目标的情况下，有权调整具体的执行步骤。
\end{enumerate}

\textbf{执行时长}

考虑到币价波动，建议不超过一季度。复杂度高的项目建议拆分阶段进行。

\textbf{提前终止}

如果开发没有按时完成，提案立即提前宣告失败。社区预留资金不会发放。

\subsection{公测}
\textbf{公测步骤}

完整的公测需要经过以下3个步骤：
\begin{enumerate}
\item 提交项目成果。开始公开测试。
\item 投票。二次投票方法和时长与一次投票相同，NAT占比y\%应大于等于x\%。
\item 技术委员会审核。通过二次投票的技术类提案，会随机派给三名技术委员会相关技术人员进行复审，全部通过即视为通过，进入发布流程。
\end{enumerate}

\textbf{测试时长}

技术类项目按照开发时长折算，最短时间为开发周期的10\%。

\subsection{发布}
非技术项目通过公测后即视为成功。技术项目需要进行链上数据更新。

\textbf{链上数据更新流程}

当协议需要升级之时，包含升级的代码会被包含在一个特殊的交易当中。这类交易只能由星云理事会暂时代为控制的主节点发起，并广播到全网。之后所有节点一旦接收到这类交易，NBRE将会自动执行其中包含的升级代码，完成更新。与此同时，核心协议的旧版本仍会在链上，以作参考。注意，节点拒绝更新只能重写本地代码。

\textbf{技术相对优势}
\begin{enumerate}
	\item 其他区块链项目升级主要还是通过硬/软分叉的方式，星云升级不需要进行分叉。
\item 交易所等不需要充停。
\item 高效率，零成本。
\end{enumerate}

举例：每个区块是火车上的一节车厢，那么每个区块中的交易就相当于车厢中的乘客。之前，如果火车发生故障，我们需要把它送到工厂维修，就像是我们手动修复漏洞，升级代码。当主网发现漏洞之后，这时，火车不能继续运行，车上的乘客也会因此受到影响。而NBRE的升级方式相当于把零部件提前放在火车上，用一节车厢专门储存零部件，当火车发生故障之时，车上的工人就可以及时换上新的部件。在此过程中，火车仍可正常运作。

\textbf{奖励发放}

发布主网后自动发放所有相关人员的奖励。


%\printbibliography
\newpage
\bibliography{reference}

\newpage 
\begin{appendices}
\section{Nebulas Autonomous Token (NAT) ​​issuance algorithm}

The issuance of NAT is based on each user's Nebulas Rank, voting behavior, and pledge amount.

\subsection{Overview}
The issuance of NAT is performed according to the weekly calculation cycle of Nebulas Rank (note: voting periods and Nebulas Rank calculations utilize the same weekly period). Based on these weekly cycles, the NAT airdrop is executed for each address on Nebulas looking at voting behavior, pledging and the previous week Nebulas Rank score.

Detailed explanation: 
For the period $i$, the new NAT $\mathcal{T}_i$ in the system is divided into three parts - the airdrop: $\mathcal{A}_i$, the voting incentive: $\mathcal{V}_i$, the pledge incentive: $\mathcal{D}_i$.
In addition, user's votes will burn/destroy a certain percentage. For example, assuming that for the cycle $i$, the NAT on the network is reduced due to voting is $\mathcal{M}_i$; then the total NAT circulation in the system is:
\begin{align}
\sum_{i=1}^{\infty} (\mathcal{A}_i + \mathcal{V}_i + \mathcal{D}_i - \mathcal{M}_i)
\end{align}

To help with understanding this section, all symbols used in this section and their corresponding explanations are listed below:
\begin{itemize}
\item $\mathcal{C}_i$: The sum of the Nebulas indices for the system in the period $i$;
\item $c_{i,j}$: User $j \in \mathcal{U}$ in the cycle $i$ Nebulas Rank Score;
\item $d_{i,j}$: User $j \in \mathcal{U}$ The total amount of NAS pledged in the period $i$;
\item $v_{i,j}$: User $j \in \mathcal{U}$ The total amount of NAT voted in the period $i$.
\end{itemize}

\subsection{Airdrop portion}
The airdrop portion is related to the user's Nebulas Rank score, specifically:
\begin{align}
    f(x) = g(x)\lambda^i
\end{align}
\noindent where $x$ is the Nebulas Rank score and can be a user's Nebulas Rank; $g(x)$ is a proportional function that adjusts the relationship between the total NAT and the total Nebulas Rank and satisfies $g(0) = 0$ ;$\lambda$ is the airdrop attenuation factor, and $\lambda < 1$.
Since $\lambda < 1$, it is easy to know $\lim_{i\to \infty}f(x) = 0$.

It can be seen that the total amount of the airdrop portion in cycle $i$ is:
\begin{align}
\mathcal{A}_i = \sum_{i=1}^{\infty}f(\mathcal{C}_i).
\end{align}

\subsection{Voting incentives portion}
The voting incentives are related to the user's voting behavior and their Nebulas Rank. For user $j \in \mathcal{U}$, the voting incentives are:
\begin{align}
\mu f(x_{i,j}) \min\{\frac{v_{i,j}}{f(x_{i-1,j})},1\}
\end{align}
\noindent where $\mu$ is the voting incentive coefficient, $\mu > 1$, which means that the user's voting behavior is given an additional reward which can be adjusted according to the number of NAS circulating in the system.

\subsection{Pledge NAS portion}
The pledge portion of NAT distribution should be less than the airdrop portion obtained based on the users improved Nebulas Rank. Based on the nature of Nebulas Rank, given the users NAS, there is an upper limit of their Nebulas Rank  score $h(d_{i,j})$~\cite{ImproveNR},
Then the NAT obtained by the pledge part is:
\begin{align}
\mathcal{D}_i = \sum_{i=1}^{\infty}\alpha f(h(d_{i,j}))
\end{align}
\noindent where $\alpha$ is the pledge incentive coefficient.


\subsection{Destroyed/Burned NAT Portion}
\label{burn}
Each time a user votes, a portion of their NAT is destroyed while the rest is returned to the user. In addition, the Nebulas Council charges a fee of $\theta\%$ for each vote in order to pay for the necessary expenses of maintaining the voting services. Therefore for each user, the destruction rate is:
\begin{align}
(1-\theta\%) \times \beta^i \times v_{i,j}
\end{align}
\noindent where $\beta$ is the destruction factor and $\beta < 1$. therefore,
\begin{align}
    \mathcal{M}_i = \sum_{i=1}^{\infty} (1-\theta\%) \times \beta^i \times v_{i,j} .
\end{align}

\subsection{Analysis}

note:
\begin{itemize}
\item The current version tentatively agrees that there is no difference between a vote and a negative vote, that is, the return ratio is different. It can then be set according to the ticket type and multiplied by a different return parameter $\mu_1$;
\item If you consider the change in the total Nebulas Rank of the system after the vote is completed, you can multiply a factor of $\mu_2$ to reflect the status of the system.
\end{itemize}

\begin{property}
	The algorithm satisfies the convergence of the total amount of NAT; in return, the total amount of NAT does not exceed the upper limit supply at any time.
\end{property}

\begin{proof}
	According to the details within the Nebulas Technical White Paper, the fixed total amount of NAS is $10^9$ with an average weekly issuance amount (on the basis of a fixed total) of $0.2\%$; in return, the total amount of NAS existing on the market for the $n$ cycle will not exceed $10^9 (1+0.002n) $
	
	Next, we must prove that the median value of assets in one cycle of all addresses (as defined in the Nebulas Rank Yellow Paper) does not exceed the total amount of existing NAS on the market. This is because for any NAS asset with a quantity of $y$, the user can only have more than half of the period (three and a half days) in one address, so it can provide the total asset value of the entire network node. Contribution of $y$.
	
	Also according to the Nebula Rank Yellow Paper, the Nebulas Rank value of any single address can not not exceed the median asset value of that address (for the same period; note that Nebulas Rank and NAT calculations are weekly and synchronized). because theNebula Rank Yellow Paper formula $Omega(CDOT)Psi(Cdot)$, the Wilbur function $Omega(CDOT)$, which is entered in the asset value, satisfies $Omega (x) leq x$, and the Access function $Psi(CDOT)$ value does not exceed 1.
	
	Combined with the above conclusions, in the $n$ cycle, the sum of all Nebulas Rank addresses does not exceed $10^9(1+0.002n)$, so that the NAT airdrop portion does not exceed $g(10^9(1+0.002n)))\lambda^n$.
	
	Also, since the NAT of the voting incentive portion does not exceed the additional portion multiplied by $\mu$, even if the increment caused by the return part is added, the total increment of NAT in the period $n$ exceeds $\mu g(10^9( 1+0.002n))\lambda^n$. In addition, the increase in the pledge portion does not exceed the total amount of NAS $g (10^9(1+0.002n))\lambda^n$.
	
	Finally, to prove the convergence of the total amount of NAT; because of the airdrop portion, the pledge portion and the incentive portion are exponentially decayed with time, it is only necessary to prove the series:
	\begin{align}
	\sum_{n=1}^{\infty} \mu g(10^9(1+0.002n))\lambda^n
	\end{align}
	Convergence, since $g(\cdot)$ is a linear function,
	\begin{align}
	\lim_{n\rightarrow \infty} \frac{\mu g(10^9(1+0.002(n+1)))\lambda^{n+1}}{\mu g(10^9(1+ 0.002n))\lambda^n} = \lambda <1
	\end{align}
	The comparison of the series can be obtained by the comparison method.
\end{proof}

The series can be convergent and verified by the comparison method.

At the same time, the above voting algorithm has the following positive properties.
\begin{enumerate}
	\item \textbf{Anti-snowball effect:} If you simply return NAT at a fixed ratio, a user can utilize all their NAT and enjoy a return greater than 1 (such as 1.1); then the total NAT will be $1.1^ The n$ has risen and the growth is too large.
	\item \textbf{Anti-buying:} If a user with a low Nebulas Rank obtains a large quantity of NAT and uses it for voting, because of the corresponding $x_{i-1}^j$ for addresses with a low Nebulas Rank score, there are very few NAT tokens returned resulting in most of them being burned. This causes the user's remaining NAT to be penalized.
	\item \textbf{Anti-inflation:} The depreciation of NAT can be effectively controlled because the distribution of additional NAT is related to the total amount of NAT in the current market.
	\item \textbf{Head-effect:} Users with a High Nebulas Rank during the early stages can have a higher total NAT quantity.
\end{enumerate}

\section{星云资产监管方案}

星云资产包括社区公共资产和星云基金会监管资产两部分。

\subsection{社区公共资产}

\subsubsection{构成}

\begin{itemize}
	\item 星云非技术白皮书资产分配相关章节中提及的社区生态预留部分：35,000,000 NAS（35\%）
    \item 共识记账系统增发，每日自动产生（ 8,219.1744 NAS/日），包括：
	    \begin{itemize}
			\item 2\%：共识记账收入 （截止节点分发之前）
			\item 1\%：星云理事会项目发展资金储备
		\end{itemize}
	\item 开发者激励协议（Developer Incentive Protocol，DIP）~\cite{mauvepaper}产生的原生激励： 自2019年5月13日起产生，初始为1\%
\end{itemize}

\subsubsection{管理}

公共资产属于社区，通过星云链上治理流程由社区共同参与管理，星云理事会负责监督。

\subsection{星云基金会资产}

\subsubsection{构成}

\begin{itemize}
	\item 星云非技术白皮书资产分配相关章节中提及的星云团队预留部分：20,000,000 NAS（20\%）
    \item 星云社区发展基金（生态投资余额）：5,000,000 NAS（5\%）
	\item 早期私募所得项目发展资金
	\item 早期生态投资所得
\end{itemize}

\subsubsection{管理}

星云基金会资产由星云基金会管理，星云基金会应确保资产整体使用情况公开透明。
\end{appendices}

\end{document}
